\documentclass[12pt]{article}

\usepackage[margin=1in]{geometry}
\usepackage[T1]{fontenc}
\usepackage{lmodern}
\usepackage{amsmath,amssymb,mathtools}
\usepackage{booktabs}
\usepackage{setspace}
\usepackage{hyperref}
\usepackage{xurl}

\hypersetup{
  colorlinks=true,
  linkcolor=black,
  citecolor=black,
  urlcolor=blue
}

\onehalfspacing
\widowpenalty=10000
\clubpenalty=10000

\title{Compton Frequency as Energy Scale and Matter-Wave Phase:\\
Expansions, Frames, and Gravitational Interference}
\author{Light-Meter Physics}
\date{}

\begin{document}

\maketitle

\begin{abstract}
Dividing a particle's energy by Planck's constant expresses that energy as a frequency scale. For a free massive particle, the rest-energy scale is the Compton frequency $\nu_C=mc^2/h$, while the energy in a chosen inertial frame gives $\nu_E=E/h=\gamma\nu_C$. This paper develops the low-speed expansion of $\nu_E$, its equivalent matter-wave dispersion relation in terms of the de Broglie wavelength, and the distinction between a plane-wave frequency at a fixed position and phase accumulated along the particle's worldline. The latter is set by the proper-time action, but this mathematical phase does not, by itself, make a particle an operational clock that emits a directly measurable signal at $\nu_C$. In a weak static gravitational field, the same rest-energy scale converts the lapse difference into a frequency-equivalent shift. The paper relates that bookkeeping to the gravitational phase in the Colella--Overhauser--Werner neutron-interferometry experiment and summarizes the contrasting interpretations of atom-interferometer tests of gravitational redshift. The central distinction is between an exact phase or energy identity and a claim about what an experiment measures.
\end{abstract}

\noindent\textbf{Keywords:} Compton frequency, matter waves, proper time, relativistic energy, atom interferometry, gravitational redshift

\section{Introduction}

The phrase ``Compton frequency'' is often used for the rest-energy scale of a massive particle expressed in cycles per second,
\begin{equation}
  \nu_C\equiv\frac{mc^2}{h}.
  \label{eq:compton-frequency}
\end{equation}
The companion angular frequency is $\omega_C=mc^2/\hbar=2\pi\nu_C$. The distinction matters: $\nu_C$ is measured in hertz, whereas $\omega_C$ is measured in radians per second. Both are derived from the particle's invariant mass and the constants $c$, $h$, and $\hbar$; neither notation alone establishes that a free particle emits a periodic signal at that rate.

This paper examines the connections among the rest-energy scale, matter-wave phase, and gravitational interference with explicit qualifications.

The paper first derives the energy and wavelength expansions, then distinguishes the relevant phase rates, and finally treats weak gravity and interferometry. The aim is to make the algebraic connections useful without confusing them with radiation, a new dynamical law, or a settled interpretation of every atom-interferometer result.

\section{Energy expressed as a frequency}

For a free particle of invariant rest mass $m>0$, the special-relativistic energy--momentum relation is
\begin{equation}
  E^2=p^2c^2+m^2c^4,
  \qquad
  E=\gamma mc^2,
  \qquad
  p=\gamma mv,
  \label{eq:energy-momentum}
\end{equation}
where $E$ and $p$ are measured in one inertial frame, $v$ is the particle's speed in that frame, and $\gamma=(1-v^2/c^2)^{-1/2}$. Define the energy-equivalent frequency
\begin{equation}
  \nu_E\equiv\frac{E}{h}.
  \label{eq:energy-frequency}
\end{equation}
Equations~\eqref{eq:compton-frequency} and \eqref{eq:energy-momentum} imply the exact relation
\begin{equation}
  \boxed{\nu_E=\gamma\nu_C.}
  \label{eq:gamma-frequency}
\end{equation}
This equation expresses frame-dependent energy in frequency units. It does not mean that the particle radiates at $\nu_E$; an emitted or detected radiation frequency requires a physical transition or signal and an appropriate measurement arrangement.

The kinetic energy $K=E-mc^2$ can likewise be expressed as a frequency-equivalent:
\begin{equation}
  \nu_K\equiv\frac{K}{h}=\nu_E-\nu_C=(\gamma-1)\nu_C.
  \label{eq:kinetic-frequency}
\end{equation}
Thus the ``beat'' $K/h$ is exactly the difference between the total-energy frequency-equivalent and the rest-energy reference. Whether that difference is an observable beat depends on whether the corresponding phases are coherently compared; it is not a generic high-frequency signal emitted by an isolated particle.

For an electron, the 2022 CODATA mass gives \cite{CODATA}
\begin{equation}
  \nu_{C,e}=\frac{m_ec^2}{h}\approx1.23559\times10^{20}\ \mathrm{Hz},
  \qquad
  \lambda_{C,e}\equiv\frac{h}{m_ec}\approx2.42631\ \mathrm{pm}.
  \label{eq:electron-values}
\end{equation}
The wavelength in Eq.~\eqref{eq:electron-values} is the (unreduced) Compton wavelength \cite{Compton}. The reduced Compton wavelength is $\hbar/(m_ec)=\lambda_{C,e}/(2\pi)$; the two conventions must not be interchanged.

\section{Low-speed expansion and a numerical example}

Let $\beta=v/c$. For $|\beta|<1$, the binomial series for $(1-\beta^2)^{-1/2}$ gives \cite{DLMF}
\begin{equation}
  \gamma
  =1+\frac{1}{2}\beta^2+\frac{3}{8}\beta^4
    +\frac{5}{16}\beta^6+\frac{35}{128}\beta^8+\cdots.
  \label{eq:gamma-series}
\end{equation}
Multiplying by $\nu_C$ yields the total-energy frequency-equivalent,
\begin{equation}
  \nu_E=\nu_C\left[
    1+\frac{1}{2}\left(\frac{v}{c}\right)^2
    +\frac{3}{8}\left(\frac{v}{c}\right)^4
    +\frac{5}{16}\left(\frac{v}{c}\right)^6
    +\frac{35}{128}\left(\frac{v}{c}\right)^8+\cdots\right].
  \label{eq:frequency-series}
\end{equation}
In energy units, the same terms are
\begin{equation}
  E=mc^2+\frac{1}{2}mv^2+\frac{3}{8}\frac{mv^4}{c^2}
    +\frac{5}{16}\frac{mv^6}{c^4}+\cdots.
  \label{eq:energy-speed-series}
\end{equation}
Consequently,
\begin{equation}
  \nu_K=\frac{mv^2}{2h}+\frac{3mv^4}{8hc^2}
    +\frac{5mv^6}{16hc^4}+\cdots.
  \label{eq:kinetic-series}
\end{equation}
The first kinetic term is the Newtonian kinetic energy divided by $h$. The series has radius of convergence $|\beta|<1$, the physical speed domain for a massive particle. A short truncation is most useful when $v/c\ll1$; convergence of the full series does not imply that a few terms are accurate near $c$.

As a worked example, take an electron at $v=0.6c$. Then $\gamma=1.25$, so Eq.~\eqref{eq:electron-values} gives the exact result $\nu_E\approx1.54449\times10^{20}\ \mathrm{Hz}$. The successive terms in the series are:
\begin{center}
\begin{tabular}{lr}
  \toprule
  Contribution & Value in $10^{20}\ \mathrm{Hz}$ \\
  \midrule
  Rest-energy baseline $\nu_C$ & 1.235590 \\
  Quadratic term & 0.222406 \\
  Quartic term & 0.060050 \\
  Sixth-order term & 0.018015 \\
  Sum through $v^6$ & 1.536061 \\
  Exact total $\gamma\nu_C$ & 1.544487 \\
  \bottomrule
\end{tabular}
\end{center}
The sixth-order truncation is about $0.546\%$ below the exact value. The missing contribution is not a physical signal omitted from the experiment; it is simply the next terms of the convergent expansion.

For a hydrogenic momentum scale $p\sim\alpha_{\mathrm{fs}} mc$, where $\alpha_{\mathrm{fs}}$ is the fine-structure constant, the first relativistic correction to the free-particle kinetic energy is smaller than its Newtonian term by order $\alpha_{\mathrm{fs}}^2$. This scaling is relevant to atomic estimates, but it is not a derivation of the full hydrogen fine structure: that requires bound-state dynamics and, for a complete result, spin-dependent and other corrections.

\section{The frequency triangle and de Broglie wavelength}

For nonzero momentum define the de Broglie wavelength $\lambda_{\mathrm{dB}}=h/p$. Divide Eq.~\eqref{eq:energy-momentum} by $h^2$ to obtain
\begin{equation}
  \boxed{\nu_E^2=\nu_C^2+\left(\frac{c}{\lambda_{\mathrm{dB}}}\right)^2.}
  \label{eq:frequency-triangle}
\end{equation}
This is the energy--momentum relation expressed as a frequency triangle: the hypotenuse is $E/h$, and the other terms are $mc^2/h$ and $pc/h=c/\lambda_{\mathrm{dB}}$. It is an algebraic form of the relativistic matter-wave dispersion relation, not a triangle of three independent emitted signals \cite{deBroglie}.

With the unreduced Compton wavelength $\lambda_C=h/(mc)$,
\begin{equation}
  \nu_E=\nu_C\sqrt{1+\left(\frac{\lambda_C}{\lambda_{\mathrm{dB}}}\right)^2}.
  \label{eq:wavelength-frequency}
\end{equation}
For long de Broglie wavelengths, $\lambda_{\mathrm{dB}}>\lambda_C$, expansion in $\lambda_C/\lambda_{\mathrm{dB}}$ gives
\begin{equation}
  \nu_E=\nu_C+\frac{h}{2m\lambda_{\mathrm{dB}}^2}
    -\frac{h^3}{8m^3c^2\lambda_{\mathrm{dB}}^4}+\cdots.
  \label{eq:long-wavelength-series}
\end{equation}
For short de Broglie wavelengths, $\lambda_{\mathrm{dB}}<\lambda_C$, an inverse-ratio expansion gives
\begin{equation}
  \nu_E=\frac{c}{\lambda_{\mathrm{dB}}}
    +\frac{c\lambda_{\mathrm{dB}}}{2\lambda_C^2}
    -\frac{c\lambda_{\mathrm{dB}}^3}{8\lambda_C^4}+\cdots.
  \label{eq:short-wavelength-series}
\end{equation}
Both binomial series converge absolutely at their shared boundary $\lambda_{\mathrm{dB}}=\lambda_C$; a finite truncation is generally most efficient well inside its respective regime. In the massless limit, $\nu_C=0$ and Eq.~\eqref{eq:frequency-triangle} reduces to $\nu_E=c/\lambda$. A photon has no rest frame and no massive-particle Compton clock, but it does have a nonzero wave frequency and phase in a specified frame.

\section{Plane-wave frequency and phase along a worldline}

In flat spacetime, a positive-energy plane-wave mode may be written
\begin{equation}
  \Psi(t,\mathbf{x})\propto
  \exp\!\left[\frac{i}{\hbar}\left(\mathbf{p}\cdot\mathbf{x}-Et\right)\right].
  \label{eq:plane-wave}
\end{equation}
At fixed spatial position its phase cycles at frequency $\nu_E=E/h=\gamma\nu_C$. Along the classical particle trajectory $\mathbf{x}(t)$, however, the phase rate is
\begin{align}
  \frac{1}{2\pi}\left|\frac{\mathrm{d}}{\mathrm{d}t}
    \frac{\mathbf{p}\cdot\mathbf{x}-Et}{\hbar}\right|
  &=\frac{E-\mathbf{p}\cdot\mathbf{v}}{h} \notag\\
  &=\frac{\nu_C}{\gamma}
   =\nu_C\frac{\mathrm{d}\tau}{\mathrm{d}t}.
  \label{eq:path-phase-rate}
\end{align}
The two rates, $\gamma\nu_C$ and $\nu_C/\gamma$, refer to different ways of sampling the same phase field. They are not contradictory: one is measured at a fixed position in the frame, and the other follows the moving particle. The de Broglie phase velocity and group velocity likewise differ: $v_{\mathrm{phase}}=E/p=c^2/v$, while the wave-packet group velocity is $v$.

The covariant origin of Eq.~\eqref{eq:path-phase-rate} is the free-particle action \cite{Landau}. Along a timelike path,
\begin{equation}
  S=-mc^2\int\mathrm{d}\tau,
  \qquad
  \varphi\equiv\frac{S}{\hbar}=-2\pi\nu_C\tau.
  \label{eq:proper-time-phase}
\end{equation}
The accumulated cycle count has magnitude $\nu_C\tau$. This is a precise statement about the phase assigned to a free-particle path. Only phase differences are observable. ``Compton clock'' is therefore shorthand for a coherent phase comparison, not a directly readable oscillator at $\nu_C$ carried by a stationary mass eigenstate.

Factoring out $e^{-imc^2t/\hbar}$ leaves a nonrelativistic Schr\"odinger envelope; its residual energy is $K$, or $K/h$ in frequency units---not a separately detectable carrier.

\section{Weak gravity and a common time reference}

Consider a static, weak gravitational field with Newtonian potential per unit mass $\Phi$, choosing its zero at a specified static reference. The lapse between proper time at a stationary location and the reference coordinate time is
\begin{equation}
  \alpha(\Phi)\simeq1+\frac{\Phi}{c^2},
  \qquad |\Phi|/c^2\ll1.
  \label{eq:lapse}
\end{equation}
For a particle measured locally by a stationary observer with speed $v_{\mathrm{loc}}$, define $\gamma_{\mathrm{loc}}=(1-v_{\mathrm{loc}}^2/c^2)^{-1/2}$. In a static metric with zero shift, the exact rate bridge is
\begin{equation}
  \frac{\mathrm{d}\tau}{\mathrm{d}t_{\mathrm{ref}}}
  =\frac{\alpha}{\gamma_{\mathrm{loc}}},
  \qquad
  \delta\tau\equiv\frac{\mathrm{d}t_{\mathrm{ref}}}{\mathrm{d}\tau}
  =\frac{\gamma_{\mathrm{loc}}}{\alpha}.
  \label{eq:static-moving-rate}
\end{equation}
The local energy is $E_{\mathrm{loc}}=\gamma_{\mathrm{loc}}mc^2$, and its energy referred to the chosen static time standard is $E_{\mathrm{ref}}=\alpha E_{\mathrm{loc}}$. Hence
\begin{equation}
  \frac{E_{\mathrm{ref}}}{mc^2}
  =\alpha\gamma_{\mathrm{loc}}
  =\alpha^2\delta\tau.
  \label{eq:static-energy-rate-bridge}
\end{equation}
In flat spacetime, $E/(mc^2)=\delta\tau=\gamma$ in a chosen inertial frame; that equality does not carry over to gravitational reference energy, for which the lapse enters both definitions. Thus, to first order in $\Phi/c^2$ and $v_{\mathrm{loc}}^2/c^2$,
\begin{align}
  \frac{E_{\mathrm{ref}}}{h}
  &\simeq\nu_C\left(1+\frac{\Phi}{c^2}
                  +\frac{v_{\mathrm{loc}}^2}{2c^2}\right) \notag\\
  &=\nu_C+\frac{m\Phi}{h}
       +\frac{mv_{\mathrm{loc}}^2}{2h}.
  \label{eq:weak-field-energy-frequency}
\end{align}
This is the energy referred to one static time standard, expressed in hertz. The zero of $\Phi$ and the reference normalization are part of the definition; $m\Phi/h$ is not a separate locally emitted frequency.

It is essential to distinguish Eq.~\eqref{eq:weak-field-energy-frequency} from the phase rate along a moving particle. In the same static coordinates,
\begin{equation}
  \nu_{\mathrm{path}}=\nu_C\frac{\mathrm{d}\tau}{\mathrm{d}t_{\mathrm{ref}}}
  =\nu_C\frac{\alpha}{\gamma_{\mathrm{loc}}}
  \simeq\nu_C\left(1+\frac{\Phi}{c^2}
                    -\frac{v_{\mathrm{loc}}^2}{2c^2}\right).
  \label{eq:weak-field-path-rate}
\end{equation}
The positive kinetic term in the reference energy at a fixed location and the negative kinematic term in the path phase rate answer different questions. For a particle held stationary, $v_{\mathrm{loc}}=0$, and the Compton phase rate per reference second is $\alpha\nu_C$. Therefore two stationary locations satisfy
\begin{equation}
  \frac{\Delta\nu_{\mathrm{stationary}}}{\nu_C}
  \simeq\frac{\Phi_2-\Phi_1}{c^2}.
  \label{eq:stationary-redshift}
\end{equation}
Equation~\eqref{eq:stationary-redshift} is a clock-rate comparison relative to the same reference time. It is not an elapsed proper time and, for an ordinary atom in a single stationary internal state, it is not by itself evidence of a directly readable Compton oscillator.

Near Earth's surface, for an upward separation $\Delta z$ small compared with Earth's radius,
\begin{equation}
  \Delta\Phi\simeq g\,\Delta z,
  \qquad
  \frac{\Delta\nu_{\mathrm{stationary}}}{\nu_C}
  \simeq\frac{g\,\Delta z}{c^2}.
  \label{eq:near-earth-redshift}
\end{equation}
For $\Delta z=1\ \mathrm{m}$, the fractional offset is about $1.09\times10^{-16}$. For cesium-133, whose Compton frequency is approximately $3.0\times10^{25}\ \mathrm{Hz}$, this corresponds to a frequency-equivalent scale of about $3.3\ \mathrm{GHz}$ per meter. That large number is the rest-mass scale multiplied by the lapse difference; it is not a claim that a cesium atom radiates a GHz signal at its Compton frequency.

\section{Gravitational phase and the COW experiment}

In the weak-field, slow-motion limit, expanding the proper-time action in Eq.~\eqref{eq:proper-time-phase} gives the Lagrangian
\begin{equation}
  L\simeq-mc^2+\frac{1}{2}mv^2-m\Phi.
  \label{eq:weak-field-lagrangian}
\end{equation}
The relative quantum phase of two paths is their action difference divided by $\hbar$. For an idealized pair of equal-duration paths held at different heights with potential difference $\Delta\Phi$, the gravitational-potential contribution has magnitude
\begin{equation}
  |\Delta\varphi_g|=\frac{m|\Delta\Phi|T}{\hbar}
   =2\pi\left(\frac{m|\Delta\Phi|}{h}\right)T
   =2\pi\nu_C\frac{|\Delta\Phi|}{c^2}T,
  \label{eq:gravity-phase}
\end{equation}
where $T$ is the time spent at the separated potentials. The sign depends on path ordering and phase convention. In a uniform field, $|\Delta\Phi|\simeq g|\Delta z|$, so the potential-energy difference can be written as a frequency-equivalent $\Delta f_{\mathrm{grav}}=mg|\Delta z|/h$; the measured relative phase accumulates as $2\pi\Delta f_{\mathrm{grav}}T$ in this idealization.

Colella, Overhauser, and Werner observed a gravity-dependent phase shift by splitting and recombining neutron matter waves \cite{COW}. In the simple geometry of horizontal arms with separation $\Delta z$ and length $L$, traversed at speed $v$, Eq.~\eqref{eq:gravity-phase} gives the familiar scale
\begin{equation}
  |\Delta\varphi_{\mathrm{COW}}|
  \simeq\frac{mgA_{\mathrm{eff}}}{\hbar v},
  \qquad A_{\mathrm{eff}}=|\Delta z|L.
  \label{eq:cow-phase}
\end{equation}
The actual interferometer phase includes the complete path and apparatus contributions. The observable is a fringe-phase difference, not direct counting of neutron Compton cycles. Expressing $mg\Delta z/h$ using $\nu_C$ is an illuminating rewrite of the same potential-energy scale, not an additional experimental result.

\section{Atom-interferometer interpretation and limits}

M\"uller, Peters, and Chu proposed that atom interferometry could be interpreted as a sensitive gravitational-redshift test using the atomic Compton frequency \cite{Muller}. Wolf and coauthors challenged that interpretation, arguing that the usual atom-interferometer observable is more properly understood as a test of free fall rather than a clock-redshift measurement \cite{Wolf}. Hohensee and coauthors published a response developing the relation between matter-wave redshift tests and equivalence-principle tests \cite{Hohensee}. These papers agree that gravity-dependent matter-wave phases are real; they disagree about how a particular interferometer's observable should be classified. The present derivations do not settle that interpretive dispute. In particular, writing a phase as $\nu_C\Delta\Phi/c^2$ does not, by algebra alone, prove that the apparatus compares two independent clocks operating at $\nu_C$.

The distinction between a phase reference and a literal carrier is also important when discussing the 2013 ``Compton clock'' demonstration by Lan and colleagues \cite{Lan}. Their experiment used an optical frequency comb to self-reference a Ramsey--Bord\'e atom interferometer and synchronize an oscillator at a technically accessible subharmonic of the Compton angular frequency. It demonstrated a clock architecture linking time measurement to particle mass; it did not directly count cycles of a $10^{25}\ \mathrm{Hz}$ oscillator with an electronic detector.

The algebraic results in this paper have a narrower scope than some clock language suggests:
\begin{itemize}
  \item $\nu_E=E/h$ is frame-dependent energy expressed as cycles per second; it is not generally a radiated frequency.
  \item $\nu_C/\gamma$ is the free-particle phase rate along its trajectory per inertial-frame second, while $\gamma\nu_C$ is the plane-wave frequency at a fixed position.
  \item The proper-time action phase is exact for a free massive particle, but a measurable clock requires an observable phase comparison.
  \item Weak-field potential terms are defined relative to a chosen static time reference; in a moving-particle calculation, energy frequency and path-phase rate must not be conflated.
  \item A matter-wave phase shift can be expressed using a Compton-frequency scale without that rewriting alone determining whether a given experiment is a redshift test, a free-fall test, or both under a specified theoretical framework.
\end{itemize}

\section{Conclusion}

The identity $\nu_E=\gamma\nu_C$ and its binomial expansion provide a compact way to express relativistic energy in frequency units. The same energy--momentum relation yields a frequency triangle in terms of the de Broglie wavelength. Separately, the free-particle action gives a phase proportional to proper time, with magnitude $2\pi\nu_C\tau$. These facts support useful carrier-and-beat language, provided the words are not mistaken for a claim that a stationary mass eigenstate emits a directly observable Compton-frequency signal.

In a static weak field, the lapse shifts stationary clock rates by $\Delta\Phi/c^2$, and the rest-energy frequency scale expresses that offset as a frequency-equivalent. Matter-wave interferometers detect relative phases that depend on gravitational potentials and path dynamics. Their interpretation as tests of Compton-frequency gravitational redshift has been debated in the literature and requires an operational and theoretical analysis beyond algebraic equivalence. The defensible conclusion is therefore precise but limited: the Compton frequency is a useful rest-energy phase scale, and its connection to proper time and gravitational phase is real; whether it functions as an experimentally tested clock depends on the observable being compared.

\begin{thebibliography}{99}
\small
\setstretch{1.0}
\raggedright
\urlstyle{same}

\bibitem{DLMF}
F. W. J. Olver, A. B. Olde Daalhuis, D. W. Lozier, B. I. Schneider, R. F. Boisvert, C. W. Clark, B. R. Miller, B. V. Saunders, H. S. Cohl, and M. A. McClain, eds.,
\emph{NIST Digital Library of Mathematical Functions}, Section 4.6, Eq.~4.6.7, National Institute of Standards and Technology,
\url{https://dlmf.nist.gov/4.6}.

\bibitem{CODATA}
P. J. Mohr, D. B. Newell, B. N. Taylor, and E. Tiesinga,
``CODATA recommended values of the fundamental physical constants: 2022,''
\emph{Reviews of Modern Physics} \textbf{97}, 025002 (2025),
\href{https://doi.org/10.1103/RevModPhys.97.025002}{doi:\nolinkurl{10.1103/RevModPhys.97.025002}}.

\bibitem{Compton}
A. H. Compton,
``A quantum theory of the scattering of X-rays by light elements,''
\emph{Physical Review} \textbf{21}, 483--502 (1923),
\href{https://doi.org/10.1103/PhysRev.21.483}{doi:\nolinkurl{10.1103/PhysRev.21.483}}.

\bibitem{deBroglie}
L. de Broglie,
``Recherches sur la th\'eorie des quanta,''
\emph{Annales de Physique} (10th series) \textbf{3}, 22--128 (1925).

\bibitem{Landau}
L. D. Landau and E. M. Lifshitz,
\emph{The Classical Theory of Fields}, 4th ed., Section 8 (Butterworth-Heinemann, Oxford, 1975).

\bibitem{COW}
R. Colella, A. W. Overhauser, and S. A. Werner,
``Observation of gravitationally induced quantum interference,''
\emph{Physical Review Letters} \textbf{34}, 1472--1474 (1975),
\href{https://doi.org/10.1103/PhysRevLett.34.1472}{doi:\nolinkurl{10.1103/PhysRevLett.34.1472}}.

\bibitem{Muller}
H. M\"uller, A. Peters, and S. Chu,
``A precision measurement of the gravitational redshift by the interference of matter waves,''
\emph{Nature} \textbf{463}, 926--929 (2010),
\href{https://doi.org/10.1038/nature08776}{doi:\nolinkurl{10.1038/nature08776}}.

\bibitem{Wolf}
P. Wolf, L. Blanchet, C. J. Bord\'e, S. Reynaud, C. Salomon, and C. Cohen-Tannoudji,
``Does an atom interferometer test the gravitational redshift at the Compton frequency?''
\emph{Classical and Quantum Gravity} \textbf{28}, 145017 (2011),
\href{https://doi.org/10.1088/0264-9381/28/14/145017}{doi:\nolinkurl{10.1088/0264-9381/28/14/145017}}.

\bibitem{Hohensee}
M. A. Hohensee, S. Chu, A. Peters, and H. M\"uller,
``Equivalence principle and gravitational redshift,''
\emph{Physical Review Letters} \textbf{106}, 151102 (2011),
\href{https://doi.org/10.1103/PhysRevLett.106.151102}{doi:\nolinkurl{10.1103/PhysRevLett.106.151102}}.

\bibitem{Lan}
S.-Y. Lan, P.-C. Kuan, B. Estey, D. English, J. M. Brown, M. A. Hohensee, and H. M\"uller,
``A clock directly linking time to a particle's mass,''
\emph{Science} \textbf{339}(6119), 554--557 (2013),
\href{https://doi.org/10.1126/science.1230767}{doi:\nolinkurl{10.1126/science.1230767}}.

\end{thebibliography}

\end{document}